\documentclass[12pt,man]{apa7}
\usepackage{styles/cwilliams-apa}

\title{Machine Learning Term Project Report}
\shorttitle{Car Crash ML Project}

\author{Christopher Williams}
\affiliation{Virginia Tech}

\course{5805: Machine Learning I}
\professor{Dr.\ Reza Jefari}
\duedate{December 4, 2025}

\begin{document}

\maketitle

\tableofcontents

\listoffigures
\listoftables

\section{Abstract}

Traffic accidents in the modern day are one of the biggest contributors to medical bills and deaths in the world. Predicting 
the severity of these accidents and their resultant injuries to develop effective road safety measures could lead to a reduction in
casualties. This paper shows the results of using machine learning models, including Logistic Regression, Random Forest, Nueral Network,
SVM, and more, to predict the severity of injury from a car crash based on certain factors. Utilizing data across 12 years (2013 - 2025, N = 209,275)
from various regions to develop and compare several models performance. The results show that Logistic Regression returned an $R^2$ of 0.105 
with the target variable total injuries whereas the Random Forest model achieved an $R^2$ of 0.54 and an F1 of 0.60. During Phase I of the
project, through a Random Forest analysis showed that the most important features were Is Rush Hour, First Crash Type: Pedestrian, Number of units (number of cars invovled), 
Lighting Condition: Daylight, and Roadway Surface Condition: Wet to be the most important features that explained the greatest amount of the variance. 
The Random Forest model did not perform as well as hoped due to a lack of helpful information in the dataset like Age of vehicle, Age of driver, impact speed, 
or Make of Vehicle.

\section{Introduction}

\section{Description of the Dataset}

\section{Phase I}

\section{Phase II}

\section{Phase III}

\section{Phase IV}

\section{Recommendations}

\section{Appendix}

\section{References}

\end{document}